\documentclass[11pt,a4paper]{article}
\usepackage[margin=2.4cm]{geometry}
\usepackage{amsmath,amssymb,amsthm}
\usepackage{graphicx}
\usepackage{booktabs}
\usepackage{array}
\usepackage{multirow}
\usepackage{longtable}
\usepackage{xcolor}
\usepackage[colorlinks=true,linkcolor=blue!50!black,urlcolor=blue!50!black,
            citecolor=blue!50!black]{hyperref}
\usepackage{caption}
\usepackage{enumitem}
\usepackage{fancyhdr}

\captionsetup{font=small,labelfont=bf}
\setlist{itemsep=1pt,topsep=3pt,parsep=0pt}

\newtheorem{lemma}{Lemma}
\newtheorem{proposition}{Proposition}

\definecolor{estab}{RGB}{20,100,40}
\definecolor{obsv}{RGB}{180,110,0}
\definecolor{conj}{RGB}{150,40,40}
\newcommand{\Est}{\textcolor{estab}{\textbf{[ESTABLISHED]}}}
\newcommand{\Obs}{\textcolor{obsv}{\textbf{[STRONG OBSERVATION]}}}
\newcommand{\Conj}{\textcolor{conj}{\textbf{[CONJECTURE]}}}
\newcommand{\Unc}{\textcolor{conj}{\textbf{[UNCERTAIN]}}}
\newcommand{\NEW}{\textcolor{blue!55!black}{\textsf{\small[new, this document]}}}

\newcommand{\fig}[1]{numerics/results/figures/final/#1}

\pagestyle{fancy}
\fancyhf{}
\fancyhead[L]{\small Dipole-conserving chain --- research progress}
\fancyhead[R]{\small 14 September 2026}
\fancyfoot[C]{\small\thepage}

\title{\vspace{-1.2cm}\textbf{Frozen sites and active-region statistics in a
dipole-conserving lattice chain}\\[2mm]
\large Research-progress record}
\author{Project: \texttt{dipole-sep-hydrodynamics} \\
\small Compiled from the repository state at commit \texttt{03b6349} plus
subsequent uncommitted work}
\date{14 September 2026}

\begin{document}
\maketitle
\thispagestyle{fancy}

\begin{abstract}
\noindent This document reconstructs the current state of the project from
the repository: source code, configuration files, raw data, notebooks, the
research log, the hypothesis file and the literature notes. Numbers quoted
here were recomputed from the raw data rather than copied from the log, and
all figures are regenerated from raw data by
\texttt{numerics/scripts/make\_progress\_figures.py}. Three verification
measurements that the project had listed as open were carried out while
preparing this document and are marked \NEW\ throughout; they are also
recorded in \texttt{RESEARCH\_LOG.md}. Interpretations are labelled
explicitly and separated from measurements, and Section~\ref{sec:problems}
is a deliberately adversarial assessment of what is weak, unsupported or
internally inconsistent in the work so far.
\end{abstract}

\section*{Executive summary}
\addcontentsline{toc}{section}{Executive summary}

\begin{itemize}
\item \textbf{Model.} A periodic chain of $L$ sites, $s_i=\pm1$ at density
exactly $1/2$. The only allowed move flips a 4-site window between
$(+,-,-,+)$ and $(-,+,+,-)$ (i.e.\ $1001\leftrightarrow0110$), conserving
both magnetisation and dipole moment. One uniformly chosen active window
flips per step.

\item \textbf{H1 (frozen fraction) --- saturation is well supported.}
Across $N=100$ to $3\times10^5$ the frozen fraction is flat.
A constant fits; logarithmic decay is excluded by a huge margin
($\chi^2_\nu\approx2.8\times10^3$); a pure power law fits only with an
exponent consistent with zero, $\alpha=-0.00083\pm0.00050$, 95\%~CI
$[-0.0018,+0.00015]$. Best estimates:
$f=0.35988\pm0.00015$ (definition A, 31 points) and
$f=0.36303\pm0.00065$ (definition B at $W=5L$, 24 points).

\item \textbf{The $W$-dependence worry is now bounded.} \NEW\ Definition B
makes the window $W$ part of the observable. Measuring $f_B(W)$ along
single trajectories shows it \emph{saturates} for $W\gtrsim20L$: the
production choice $W=5L$ overshoots the plateau by only
$+0.003$ to $+0.004$, and $W=10L$ by $\lesssim0.001$.

\item \textbf{The estimator is validated against ground truth.} \NEW\ Exact
breadth-first enumeration of the Krylov sector at $L\le60$ gives the
\emph{static} frozen fraction that H1 actually defines. Once
fully-jammed configurations --- which the dynamical estimator structurally
cannot sample, and which vanish exponentially with $L$ --- are excluded from
both sides, exact and dynamical agree with $\chi^2_\nu=0.83$ over six
matched $L$. This closes the \Unc\ flag H1 has carried since it was stated.

\item \textbf{Strong fragmentation, measured directly.} \NEW\ Median Krylov
sector size grows as $(1.135)^L$ against $(1.969)^L$ half-filled
configurations, so a sector occupies a fraction $\sim(0.576)^L$ of its
symmetry sector.

\item \textbf{H2 (active-region lengths) --- exponential, not power law.}
The distribution of active-region lengths is exponential with
$\lambda\approx13$ sites, \emph{independent of $L$} over
$L=100$--$12800$. A full-range exponential fit beats a power law at all
eight $L$ ($R^2=0.910$--$0.966$ versus $0.872$--$0.906$), and the rescaled
distributions collapse. Mean nonzero length $15.4$--$16.1$ and zero-gap
rate $0.885$--$0.889$ are both flat in $L$.

\item \textbf{The forbidden lengths are now explained.} \NEW\ Active-region
lengths $1,2,3$ and $5$ never occur; $4$ and all $\ell\ge6$ do. An
exhaustive finite check (zero violations over all $2^5$ configurations)
plus a short argument shows this is exact, not statistical: two window
positions one site apart can never both fire in isolation. \emph{This
derivation has not yet been reviewed by the researcher.}

\item \textbf{Physical reading (interpretation, not measurement).} A finite
frozen fraction with exponentially distributed, $L$-independent active
regions means active clusters stay microscopic as $L\to\infty$. In the
language of the project README, this is the ``small clusters'' branch:
there is no scale-free active structure available to carry an anomalous
hydrodynamic tail.

\item \textbf{Main open weaknesses.} No production run above $N=10^4$ under
the current definition; \texttt{HYPOTHESES.md} still states a static order
parameter while every production number is a dynamical proxy; the
$W$-scan settings comparison changed four parameters at once; and the
literature's exact analytic predictions have not been evaluated
numerically to compare against our $0.36$.
\end{itemize}

\tableofcontents
\newpage

\section{Research question and motivation}

\subsection{The model}
The system is a ring of $L$ sites carrying $s_i\in\{+1,-1\}$ (equivalently
occupation $n_i\in\{1,0\}$), at density exactly $1/2$. Write a 4-site window
starting at $i$ as $(s_i,s_{i+1},s_{i+2},s_{i+3})$ with indices mod $L$. The
window is \emph{active} iff it equals
\begin{equation}
A=(+1,-1,-1,+1) \quad\text{or}\quad B=(-1,+1,+1,-1),
\end{equation}
i.e.\ $1001$ or $0110$ in occupation language, and the only allowed move is
$A\leftrightarrow B$ on an active window. The dynamics is a classical
continuous-time-like stochastic process implemented as: at each step, pick
one active window uniformly at random and flip it
(\texttt{numerics/src/dipole\_chain.py}).

The move conserves the particle number $Q=\sum_i n_i$ and the dipole moment
$P=\sum_i i\,n_i$ (modulo $L$ under periodic boundaries), because it moves
two particles in opposite directions by equal amounts. This is the defining
structure of a dipole-conserving / ``fracton'' lattice gas: single particles
are immobile, only correlated pairs move.

A useful exact fact, used below: for a uniformly random half-filled
configuration the probability that a given window is active is
$2/2^4=1/8$, so $\langle n_{\text{active}}\rangle=L/8=12.5\%$ of $L$ at
$t=0$; the research log records that this stays at $12.5$--$13.6\%$ in
steady state across $L=10^2$--$10^5$.

\subsection{Why the problem is interesting}
The project README states the motivation directly. Because the kinetic
constraint is strong, ``some regions do not evolve at all''. The programme
is then:
\begin{enumerate}
\item establish that a \emph{finite fraction} of the chain is frozen in the
thermodynamic limit (not a vanishing one);
\item determine whether the length distribution of the \emph{active} regions
between frozen sites is exponential, $\sim e^{-\ell/\eta}$, or a power law,
$\sim \ell^{-\alpha}$.
\end{enumerate}
The second question is the one with physical consequences. If the
distribution is exponential, active regions are microscopic clusters whose
size does not grow with $L$. If it is a power law, arbitrarily large active
clusters survive in the thermodynamic limit, and the README's target
observable --- the global density--density correlation function --- could
acquire an anomalous tail: even with a finite frozen fraction arresting
global transport, a perturbation could still relax globally through rare
large active regions. The README explicitly contemplates abandoning this
model for another if the distribution turns out exponential.

\subsection{Hypotheses as stated in the project}
\begin{description}[style=nextline,leftmargin=1.2cm]
\item[H1 (\texttt{HYPOTHESES.md}, 12 Sep)] $f_N\to f_\infty>0$ as
$N\to\infty$ (candidate H1a), against H1b $f_N\sim cN^{-\alpha}$ and
H1c $f_N\sim c/\ln N$. The stated order parameter is \emph{static}: a site
is frozen iff no sequence of allowed moves starting from the initial
configuration can ever change it. The file carries a flagged
\Unc: this static/dynamical equivalence had never been verified by
enumeration.
\item[H2 (14 Sep)] The active-region length distribution either collapses
onto an $L$-independent memoryless form (H2a) or retains an anomalously
heavy tail (H2b). The measured quantity is the gap between consecutive
frozen sites at a single post-burn-in snapshot, with zero-length gaps
discarded by explicit researcher instruction.
\item[Q3 (\texttt{QUESTIONS.md} \#3, 14 Sep)] Why do active-region lengths
$1,2,3,5$ never occur while $4$ and all $\ell\ge6$ do?
\end{description}

\subsection{Literature context contained in the project}
Seven papers were read and assessed (\texttt{LITERATURE.md}; applicable
details in \texttt{numerics/theory/Notes/}). The directly relevant content:
\begin{itemize}
\item \textbf{Classen-Howes (2024) thesis, Ch.~5 §5.1} treats exactly this
model ($d=2$ hardcore, range $k=4$). Its general result is a
strong-to-weak fragmentation transition at critical filling
$\nu_c=1/(k-2)$; for $k=4$ this is $\nu_c=1/2$, and $(d,k)=(2,4)$ is one of
only two degenerate cases where $\nu_c$ coincides with its particle--hole
mirror, so the weakly fragmented phase disappears and the model is
\emph{strongly fragmented at every filling}. It also derives an exact
active-bubble density $\propto\nu^x(1-\nu)^{\ell-x}$, exponential in
$\ell$, and gives $\ell_{\min}=\ell_{\max}=4$ for the minimal two-particle
bubble.
\item \textbf{Morningstar, Khemani \& Huse (2020)} study the $n_i\in\{0,1,2\}$
version; our model is exactly its hardcore invariant sub-sector, and it uses
the same static Krylov-based definition of a frozen site.
\item \textbf{Moudgalya \& Motrunich (2021)} supplies a reusable exact
counting technique (canonical form $\to$ tiling $\to$ linear recursion).
\item Three further papers (Han--Lake--Ro; Glorioso \textit{et al.};
Zerba \textit{et al.}) treat the complementary dense/ergodic regime and do
not bear on H1/H2. Sala \textit{et al.} (2020), flagged by three of them as
the likely key reference, was checked and found \emph{not} to apply: it is a
spin-1 model with a vacancy, and blocks of two like charges --- the interior
of our only active pattern --- are frozen there.
\end{itemize}

\section{Work completed}

\subsection{Chronology}
\begin{center}\small
\begin{tabular}{@{}llp{9.4cm}@{}}
\toprule
Date & Stage & Content \\
\midrule
12 Sep & H1 v1 & Pipeline built; definition A (cumulative) runs to
$N=102400$; fits (constant vs.\ power-law correction). \\
13 Sep & H1 v1 & \texttt{fitting\_range} run to $N=3\times10^5$ (11h29m);
combined 31-point fits; H1b/H1c added; overshoot/commensuration checks. \\
13 Sep & \textbf{redefinition} & ``Frozen'' redefined as a rolling window
(definition B), $W=5L$, $M=10$, on researcher instruction. \\
13 Sep & validation & Fixed-patience burn-in found to fail at large $N$
(2/3 trials drifting at $N=10^5$); two fixes compared; self-scaling
patience (``method 1'') locked in. \\
13 Sep & H1 v2 & \texttt{production\_intermediate} $N=100$--$10^4$ (1h28m);
re-run with doubled $W,M,c,n_{\min}$ (6h03m); settings comparison. \\
14 Sep & H2 & Gap-length pipeline built; test (11s) and production
($L=100$--$12800$, 14m22s) runs; analysis notebook. \\
14 Sep & literature & Seven papers reconstructed; \texttt{LITERATURE.md}
and \texttt{Notes/} written. \\
14 Sep & \NEW & $f_B(W)$ scan; exact BFS static frozen fraction;
small-$L$ like-for-like validation; forbidden-length argument. \\
\bottomrule
\end{tabular}
\end{center}

\subsection{Three inequivalent definitions of ``frozen''}
Much of the project's subtlety lives here, and the three definitions are
easy to conflate.

\begin{center}\small
\begin{tabular}{@{}lp{6.0cm}p{6.2cm}@{}}
\toprule
& \textbf{Definition} & \textbf{Status in the project} \\
\midrule
\textbf{Static} &
Site takes the same value in \emph{every} configuration reachable by any
move sequence. A property of the initial string only. &
What \texttt{HYPOTHESES.md} H1 literally states. Never measured until this
document; now measured exactly at $L\le60$ by BFS. \\
\addlinespace
\textbf{A (cumulative)} &
Site has never flipped since $t=0$; run stopped once no \emph{new} site has
flipped for $20L$ steps. &
All 12--13 Sep production data ($N\le3\times10^5$). Superseded, archived
(not deleted). \\
\addlinespace
\textbf{B (rolling)} &
Site has not flipped during the trailing window $[t-W,t]$; sites may
re-enter the frozen set. &
Current production definition; all H1 v2 and all H2 data. \\
\bottomrule
\end{tabular}
\end{center}

Two structural facts follow immediately and are worth stating because they
constrain everything:
\begin{equation}
f_{\text{static}} \;\le\; f_B(W') \;\le\; f_B(W)
\quad\text{for } W\le W' ,
\end{equation}
pointwise on any trajectory, since a longer window gives a site more chances
to disqualify itself; and definition A is definition B evaluated over the
window $[0,T_A]$, so it is directly comparable once $T_A$ is known
(measured below: $T_A\approx27$--$32\,L$).

\textbf{$W$ is part of the observable, not the stopping rule.} Every H1 v2
number is $f_B$ at one particular $W$. This was flagged when definition B
was adopted but never quantified until Section~\ref{sec:wdep}.

\subsection{Code and computational method}
\begin{itemize}
\item \texttt{src/dipole\_chain.py} --- the chain, with an $O(1)$
active-site list (swap-with-last) and a local rescan of the ten windows
that a flip can affect.
\item \texttt{src/frozen\_fraction.py} --- definition A
(\texttt{measure\_frozen\_fraction}), definition B
(\texttt{measure\_frozen\_fraction\_rolling}), the shared self-scaling
burn-in \texttt{\_self\_scaling\_burnin}, adaptive ensemble drivers.
\item \texttt{src/burnin\_comparison.py} --- kept as the record of the
method-1 vs.\ method-2 comparison.
\item \texttt{src/gap\_length.py} --- H2 measurement; cyclic gap extraction
from the site-level frozen mask.
\item \NEW\ \texttt{src/w\_dependence.py}, \texttt{src/static\_frozen.py}
--- the two new measurements, with matching scripts and configs.
\item Ensembles use \texttt{numpy.random.SeedSequence} spawning with a
per-estimator integer tag (1 rolling, 2 gap-length, 3 static, 4 $W$-scan),
so every estimator has an independent, reproducible stream.
\item Sampling is adaptive: draw batches until the relative standard error
on the mean falls below a target (1\% in production) subject to a
\texttt{min\_samples} floor and a cap.
\end{itemize}

\subsection{The burn-in problem, and why it matters}
This was the most consequential methodological episode in the project. The
original burn-in used a fixed patience of $20L$ steps. The measured
growth-phase duration in units of $L$ grows with $N$, so the ratio of
patience to growth time fell from ${\sim}5\times$ at $N=10^2$ to
${\sim}0.4\times$ at $N=10^6$. A trend test across the $M$ measurement
windows then found significant drift in 1/5 trials at $N=3\times10^4$ and
2/3 at $N=10^5$ --- i.e.\ measurements were starting before the system had
settled, exactly where the ratio argument predicted.

The adopted fix (``method 1'') extends burn-in in $W$-length windows until
$(t-G)\ge c\,G$, where $G$ is the step of the most recent first-time flip;
this pins the patience/growth ratio at $c$ for every $N$. It gave 0/10
significant trends at the same $N$, and agreed with an independent
drift-gated scheme to within $\pm0.0004$ on identical trajectories. A
zero-variance \texttt{NaN} bug in the trend test (which had been silently
counted as ``not significant'') was fixed in the same pass.

\subsection{What is completed, ongoing, and untested}
\begin{description}[style=nextline,leftmargin=0.9cm]
\item[Completed] H1 under definition A up to $N=3\times10^5$ with full
model comparison; H1 under definition B up to $N=10^4$ at two settings; H2
production $L=100$--$12800$ with shape/tail analysis; burn-in validation;
estimator validation against exact enumeration; $f_B(W)$ scan; literature
review.
\item[Ongoing / configured but not run] \texttt{production\_large\_scale.yaml}
($N=10^4$--$2\times10^4$, estimated 4.2--4.6\,h) is configured and was
approved, but a scheduled launch was interrupted and it has never run. A
reverted \texttt{production\_intermediate} configuration ($W=5L$,
$n_{\min}=40$) was costed at ${\sim}2.2$\,h and also never run.
\item[Proposed, never attempted] Extending the Classen-Howes Appendix~B.6
recursion to an exact $\rho_F(\nu=1/2)$ and an exact multiplicity table; the
paired $W=5L$ vs.\ $W=10L$ disentangling test on identical trajectories
(now largely superseded by the $W$-scan); density--density correlation
functions, which are the README's actual long-term target and have not been
touched.
\end{description}

\section{Results}
\label{sec:results}

\subsection{H1: the frozen fraction is flat in $N$}

\begin{figure}[htbp]
\centering
\includegraphics[width=\textwidth]{\fig{fig_h1_overview.pdf}}
\caption{All frozen-fraction data in the project. (a) Definition A
(archived, 31 points, $N=100$--$3\times10^5$), definition B at $W=5L$
(24 points) and at $W=10L$ (20 points); dashed lines are the
inverse-variance weighted mean of each series. (b) The same data as
residuals from each series' own constant fit. The two high points with
large error bars near $N=200$--$400$ come from the \texttt{test.yaml} smoke
run, whose target relative error was 10\% rather than 1\%. Error bars are
$\pm1$ standard error of the ensemble mean throughout.}
\label{fig:h1overview}
\end{figure}

Table~\ref{tab:h1fits} collects the model comparison, recomputed here from
the raw \texttt{.npy} samples rather than copied from the log. The
conventions are those already used in the project: weighted least squares
with weights $1/\sigma^2$, $\chi^2=\sum((f-\text{model})/\sigma)^2$.

\begin{table}[htbp]
\centering\small
\caption{Model comparison for H1. ``$\alpha$'' in H1b is the decay exponent,
so $\alpha>0$ means $f_N$ vanishes. Negative fitted $\alpha$ means the best
fit is a slight \emph{increase} with $N$.}
\label{tab:h1fits}
\begin{tabular}{@{}llrrr@{}}
\toprule
Dataset & Model & parameters & $\chi^2/\text{dof}$ & $\chi^2_\nu$\\
\midrule
\multirow{4}{*}{\parbox{4.3cm}{Definition A, 31 pts,\\ $N=100$--$3\times10^5$}}
 & H1a constant & $f_\infty=0.35988\pm0.00015$ & $47.16/30$ & 1.572 \\
 & H1a $+aN^{-\alpha}$ & $f_\infty=0.35990\pm0.00020$, $\alpha=1.28\pm0.78$ & $33.36/28$ & 1.192 \\
 & H1b $cN^{-\alpha}$ & $\alpha=-0.00083\pm0.00050$ & $44.39/29$ & 1.531 \\
 & H1c $c/\ln N$ & $c=4.218$ & $83289/30$ & \textbf{2776} \\
\addlinespace
\multirow{4}{*}{\parbox{4.3cm}{Definition B, $W=5L$,\\ 24 pts, $N=100$--$10^4$}}
 & H1a constant & $f_\infty=0.36303\pm0.00065$ & $24.71/23$ & 1.074 \\
 & H1a $+aN^{-\alpha}$ & $f_\infty=0.36430\pm0.00240$, $\alpha=0.63\pm0.63$ & $17.41/21$ & 0.829 \\
 & H1b $cN^{-\alpha}$ & $\alpha=-0.00443\pm0.00180$ & $18.62/22$ & 0.846 \\
 & H1c $c/\ln N$ & $c=2.886$ & $9655/23$ & \textbf{420} \\
\addlinespace
\multirow{2}{*}{\parbox{4.3cm}{Definition B, $W=10L$,\\ 20 pts}}
 & H1a constant & $f_\infty=0.36043\pm0.00057$ & $28.11/19$ & 1.479 \\
 & H1c $c/\ln N$ & $c=2.935$ & $10318/19$ & \textbf{543} \\
\bottomrule
\end{tabular}
\end{table}

\begin{figure}[htbp]
\centering
\includegraphics[width=0.72\textwidth]{\fig{fig_h1_models.pdf}}
\caption{The four candidate forms fitted to the definition-A dataset (the
widest range available, $N=100$ to $3\times10^5$). H1c ($c/\ln N$, dotted)
fails catastrophically. H1b's best fit is indistinguishable from a constant
because its fitted exponent is consistent with zero. Lower panel: residuals
in units of $\sigma$.}
\label{fig:h1models}
\end{figure}

\paragraph{Reading of Table~\ref{tab:h1fits}.}
\begin{itemize}
\item \textbf{H1c is excluded} under every definition, by three orders of
magnitude in $\chi^2_\nu$. This is not a marginal call.
\item \textbf{H1b is not excluded as a functional form}, but its fitted
exponent is consistent with zero and slightly negative under all three
datasets; the 95\% confidence interval from definition A is
$\alpha\in[-0.0018,+0.00015]$. Taking the upper end at face value,
$f_N$ could fall by at most ${\sim}0.03\%$ per decade of $N$. Over the
four decades measured this is unobservable, and the fit is degenerate with
a constant.
\item \textbf{H1a fits.} For definition B at $W=5L$ the constant model has
$\chi^2_\nu=1.07$ ($p=0.37$): the data are entirely consistent with a
constant. For definition A, $\chi^2_\nu=1.57$ ($p=0.024$) --- a mild but
real excess of scatter over the quoted error bars, discussed in
Section~\ref{sec:problems}.
\end{itemize}

\subsection{The settings comparison, and why it is confounded}

\begin{figure}[htbp]
\centering
\includegraphics[width=0.68\textwidth]{\fig{fig_h1_settings.pdf}}
\caption{The 13 September re-run. Four parameters were changed
simultaneously ($W$, $M$, $c_{\text{selfscale}}$, $n_{\min}$ all doubled).
Lower panel: the shift is small, systematic and essentially
$N$-independent.}
\label{fig:h1settings}
\end{figure}

The doubled-settings re-run shifted $19$ of $20$ points downward by a
weighted mean of $-0.00298\pm0.00088$ ($3.4\sigma$ pooled; no individual
point exceeds $1.96\sigma$). The shift shows no dependence on system size
(slope against $\ln N$: $-0.00013\pm0.00036$, $p=0.73$).

Two of the four changed parameters independently predict a downward shift:
doubling $W$ must lower $f_B$ by the monotonicity above, and doubling
$c_{\text{selfscale}}$ lengthens burn-in. The other two ($M$, $n_{\min}$)
only reduce variance. As designed, the experiment cannot separate these.
The $W$-scan below settles it: the $W$ effect alone accounts for
$0.003$--$0.004$, which is the entire observed shift, leaving nothing that
requires a burn-in explanation.

\subsection{$W$-dependence of definition B}
\label{sec:wdep}

\begin{figure}[htbp]
\centering
\includegraphics[width=0.66\textwidth]{\fig{fig_wdep.pdf}}
\caption{\NEW\ $f_B$ as a function of the trailing-window length, with every
$W$ read off the \emph{same} trajectories, so the $W$-dependence carries no
realisation noise. The curve is flat to the fourth decimal from
$W\approx20L$ onward at all three $L$. Dashed line: the definition-A
weighted mean over $N\ge100$ for reference (a different estimator over a
different $N$ range, shown only to set the scale).}
\label{fig:wdep}
\end{figure}

\begin{table}[htbp]
\centering\small
\caption{\NEW\ $f_B(W)$ from a single span of $200L$ steps per realisation.
Standard errors are $0.0057$--$0.0084$; because all columns come from the
same trajectories, \emph{differences} between columns are far more precise
than these errors suggest.}
\label{tab:wdep}
\begin{tabular}{@{}lrrrrrrrr@{}}
\toprule
$W/L$ & 1 & 2 & 5 & 10 & 20 & 50 & 100 & 200\\
\midrule
$L=250$  & 0.4067 & 0.3649 & 0.3505 & 0.3483 & 0.3474 & 0.3474 & 0.3474 & 0.3474\\
$L=1000$ & 0.4187 & 0.3848 & 0.3662 & 0.3626 & 0.3625 & 0.3625 & 0.3625 & 0.3625\\
$L=4000$ & 0.4184 & 0.3828 & 0.3686 & 0.3653 & 0.3648 & 0.3647 & 0.3647 & 0.3647\\
\bottomrule
\end{tabular}
\end{table}

The production window $W=5L$ overshoots the plateau by $+0.0031$, $+0.0037$
and $+0.0039$ at $L=250,1000,4000$; $W=10L$ by $+0.0009$, $+0.0001$,
$+0.0006$. So the definitional freedom in H1 v2 costs at most about $1\%$
of the value and does not grow with $L$.

I also measured the definition-A estimator's effective observation window
directly: $T_A=26.7L$ at $L=1000$ and $32.1L$ at $L=4000$ (30 realisations
each). Both lie inside the saturated regime of $f_B(W)$, which explains why
definitions A and B give consistent answers at matched $L$ despite being
constructed differently.

\subsection{Validation against exact enumeration}
\label{sec:bfs}

\begin{figure}[htbp]
\centering
\includegraphics[width=\textwidth]{\fig{fig_static_bfs.pdf}}
\caption{\NEW\ (a) The exact static frozen fraction from breadth-first
enumeration of the Krylov sector, with and without fully-jammed
configurations, against the dynamical estimator at saturated $W$.
(b) The fraction of random half-filled configurations that are completely
jammed (Krylov sector = one state), which is what forces the two curves in
(a) apart at small $L$. (c) Median Krylov sector size against the total
number of half-filled configurations.}
\label{fig:bfs}
\end{figure}

\begin{table}[htbp]
\centering\small
\caption{\NEW\ Estimator validation. ``Exact, jammed removed'' is the
sub-ensemble the dynamical estimator can actually sample, since burn-in
reports jammed realisations as non-converged and discards them.}
\label{tab:validation}
\begin{tabular}{@{}rrrrr@{}}
\toprule
$L$ & exact $f_{\text{stat}}$ (all) & exact, jammed removed & dynamical $f_B(W{=}500L)$ & difference \\
\midrule
16 & $0.2863\pm0.0075$ & $0.2114\pm0.0060$ & $0.2132\pm0.0137$ & $+0.0019$ ($+0.1\sigma$)\\
24 & $0.3023\pm0.0060$ & $0.2815\pm0.0055$ & $0.2694\pm0.0120$ & $-0.0121$ ($-0.9\sigma$)\\
32 & $0.3086\pm0.0051$ & $0.3033\pm0.0050$ & $0.2981\pm0.0113$ & $-0.0052$ ($-0.4\sigma$)\\
40 & $0.3274\pm0.0047$ & $0.3250\pm0.0046$ & $0.3438\pm0.0127$ & $+0.0188$ ($+1.4\sigma$)\\
48 & $0.3347\pm0.0060$ & $0.3334\pm0.0059$ & $0.3171\pm0.0106$ & $-0.0163$ ($-1.3\sigma$)\\
56 & $0.3286\pm0.0085$ & $0.3286\pm0.0085$ & $0.3225\pm0.0102$ & $-0.0061$ ($-0.5\sigma$)\\
\midrule
\multicolumn{4}{r}{$\chi^2=4.97$ on 6 points} & $\chi^2_\nu=0.83$\\
\bottomrule
\end{tabular}
\end{table}

Three things come out of this.

\textbf{(i) The estimator is unbiased where it can be checked.} Against
ground truth, on the sub-ensemble it can sample, the dynamical estimator
shows no detectable bias ($\chi^2_\nu=0.83$). This is the check H1 has been
asking for since it was written.

\textbf{(ii) A real small-$L$ artifact is identified.} Fully jammed
configurations --- where no window is active at all, so every site is
trivially frozen --- make up $16.2\%$ of random half-filled configurations
at $L=12$, $9.5\%$ at $L=16$, $0.75\%$ at $L=32$ and $0\%$ by $L=56$. The
dynamical estimator discards them, which biases it low by up to $4.7\sigma$
at $L=16$. At production sizes ($L\ge100$) every run reports
\texttt{n\_discarded}$=0$, so the artifact is absent there --- but note that
this is an extrapolation of an exponentially decaying jamming probability,
not a direct check at production $L$.

\textbf{(iii) $f_{\text{static}}$ is still rising at $L=60$.} It goes from
$0.286$ at $L=16$ to $0.337\pm0.010$ at $L=60$, so the exact small-$L$
numbers cannot be compared directly with the $0.36$ measured at
$N\ge100$; they are consistent with approaching it but do not by themselves
establish the asymptote.

\paragraph{Strong fragmentation, measured.} Fitting the median sector size
gives $(1.135)^L$ against $\binom{L}{L/2}\sim(1.969)^L$, so a typical
Krylov sector occupies a fraction $\sim(0.576)^L$ of its symmetry sector.
This is a direct measurement of strong Hilbert-space fragmentation in
\emph{this} model, and it independently reproduces what the Classen-Howes
thesis predicts for $(d,k)=(2,4)$ at every filling.

\subsection{H2: active-region length distribution}

\begin{table}[htbp]
\centering\small
\caption{H2 production run
(\texttt{gap\_length\_gap\_length\_production\_20260914T045549Z}), $W=5L$,
$c_{\text{selfscale}}=5$, one post-burn-in snapshot per realisation,
zero-length gaps excluded from $P_L(\ell)$. $\lambda$ and $\alpha$ are
full-range OLS fits of $\ln P$ against $\ell$ and against $\ln\ell$.}
\label{tab:h2}
\begin{tabular}{@{}rrrrrrrrr@{}}
\toprule
$L$ & realisations & pooled gaps & $\langle\ell\rangle$ & $P(\ell{=}0)$
& $\lambda$ & $R^2_{\exp}$ & $\alpha$ & $R^2_{\text{pow}}$\\
\midrule
100   & 1400 & 5643  & 16.115 & 0.8850 & 12.95 & 0.952 & 2.42 & 0.879\\
200   & 720  & 5825  & 15.730 & 0.8888 & 13.26 & 0.917 & 2.49 & 0.872\\
400   & 320  & 5261  & 15.529 & 0.8864 & 13.27 & 0.939 & 2.54 & 0.897\\
800   & 160  & 5225  & 15.633 & 0.8872 & 13.22 & 0.943 & 2.43 & 0.895\\
1600  & 100  & 6544  & 15.554 & 0.8876 & 13.76 & 0.910 & 2.51 & 0.892\\
3200  & 60   & 7793  & 15.665 & 0.8886 & 13.08 & 0.966 & 2.60 & 0.904\\
6400  & 40   & 10416 & 15.636 & 0.8882 & 12.89 & 0.952 & 2.72 & 0.889\\
12800 & 20   & 10530 & 15.407 & 0.8877 & 14.15 & 0.915 & 2.72 & 0.901\\
\bottomrule
\end{tabular}
\end{table}

\begin{figure}[htbp]
\centering
\includegraphics[width=\textwidth]{\fig{fig_h2_pmf.pdf}}
\caption{The measured length distribution on the two diagnostic axes. (a)
Semi-log: the data are straight over more than two decades in $P$, which is
the exponential signature. (b) The same data log--log, where a power law
would be straight; it is visibly curved. The flat floor near
$P\sim10^{-4}$ in both panels is the single-count resolution limit
($1/n_{\text{pooled}}$), not physics.}
\label{fig:h2pmf}
\end{figure}

\begin{figure}[htbp]
\centering
\includegraphics[width=0.6\textwidth]{\fig{fig_h2_collapse.pdf}}
\caption{Rescaled distributions for all eight $L$. The collapse is the
H2a signature: the shape is $L$-independent, and close to the memoryless
reference $e^{-x}$.}
\label{fig:h2collapse}
\end{figure}

\begin{figure}[htbp]
\centering
\includegraphics[width=\textwidth]{\fig{fig_h2_vs_L.pdf}}
\caption{$L$-dependence of the summary statistics. All three are flat over
two decades in $L$. In (c) the power-law exponent is scaled by 4 to share
the axis.}
\label{fig:h2vsL}
\end{figure}

\paragraph{What the numbers say.} The exponential form is preferred over a
power law at \emph{every} $L$ when both are fitted over the full range
($R^2_{\exp}=0.910$--$0.966$ against $R^2_{\text{pow}}=0.872$--$0.906$).
The decay length $\lambda\approx12.9$--$14.2$ sites shows no systematic
trend with $L$ over two decades, as do
$\langle\ell\rangle=15.4$--$16.1$ and the zero-gap rate
$0.885$--$0.889$.

A caveat the project already recorded: if the fit is restricted to the tail
($\ell\ge15$) the two forms become nearly indistinguishable
($R^2_{\exp}=0.876$--$0.954$ vs.\ $R^2_{\text{pow}}=0.897$--$0.951$, with
the power law marginally ahead at five of eight $L$). The full-range
comparison is the more discriminating one, because it includes the region
around the peak where a power law would have to diverge and does not; but
it is worth being explicit that the \emph{far tail alone} does not by itself
settle the question, and no likelihood-based model selection (AIC, or a
likelihood-ratio test on the untruncated samples) has been done.

\subsection{Consistency checks}
Four independent checks were run while preparing this document; all pass.

\begin{enumerate}
\item \textbf{The dynamics conserves what it claims to.} \NEW\ Over
$1.44\times10^5$ steps spanning $L=40,100,400$ and 12 seeds each:
magnetisation $\sum_i s_i$ was conserved exactly (0 failures), the dipole
moment was conserved exactly modulo $L$ (0 failures), and the incremental
$O(1)$ active-site list agreed with a full brute-force rescan at every one
of 720 spot checks. The unwrapped dipole was observed to take only the
values $\{0,-L\}$ relative to its initial value, i.e.\ it is conserved
exactly except for the expected jump of $L$ when a move wraps the periodic
boundary. This had never been tested in the repository.
\item \textbf{Cross-pipeline.} The H2 gap-length run records the snapshot
frozen fraction as a by-product, computed by different code from the H1
estimator (single window rather than $M$ averaged windows, separate seed
stream). At the five overlapping $L$ the two agree within
$-1.1\sigma$ to $+1.0\sigma$.
\item \textbf{Exact identity.} Since every frozen site opens exactly one
gap on the ring, the mean over \emph{all} gaps (zeros included) must equal
$(1-f)/f$. Recomputing from the stored counts reproduces this to a ratio of
$1.0000$ at every $L$, confirming the gap extraction and the bookkeeping of
discarded zeros are internally consistent.
\item \textbf{Estimator vs.\ ground truth.} Table~\ref{tab:validation}.
\end{enumerate}

\subsection{The forbidden lengths: an exact argument}
\label{sec:forbidden}

\begin{figure}[htbp]
\centering
\includegraphics[width=0.72\textwidth]{\fig{fig_h2_forbidden.pdf}}
\caption{Pooled counts at small $\ell$. Shaded columns have \emph{exactly
zero} counts at every $L$, across ${\sim}57{,}000$ pooled gaps. $\ell=0$ is
excluded by construction; $\ell=1,2,3,5$ are absent dynamically.}
\label{fig:forbidden}
\end{figure}

Let $S\subseteq\mathbb{Z}_L$ be the set of window positions that fired at
least once during the measurement window, and $F=\bigcup_{p\in S}\{p,p{+}1,
p{+}2,p{+}3\}$ the set of sites recorded as active. A maximal run of $F$
whose fired positions run from $p_1$ to $p_k$ (consecutive gaps $\le4$, so
the 4-blocks chain) covers $[p_1,p_k+3]$, hence has length
\begin{equation}
\ell=p_k-p_1+4 .
\end{equation}
So $\ell\ge4$ always, excluding $\ell\in\{1,2,3\}$, and $\ell=5$ requires
$p_k-p_1=1$: exactly the two positions $p$ and $p+1$ fired, and nothing else
in that component.

\begin{lemma}[verified exhaustively]
\label{lem:offset1}
For every configuration of five consecutive sites:
(i) window positions $p$ and $p+1$ are never simultaneously active;
(ii) firing $p$ leaves $p+1$ inactive; (iii) firing $p+1$ leaves $p$
inactive.
\end{lemma}
\begin{proof}[Verification]
Direct enumeration of all $2^5=32$ configurations gives $0$ violations of
each of (i), (ii), (iii). For contrast, offsets $2,3,4$ admit $2,2,4$
simultaneously-active configurations respectively, which is what makes
$\ell=6,7,8$ reachable.
\end{proof}

\begin{proposition}
\label{prop:five}
Under this dynamics, a maximal active run of length exactly $5$ cannot
occur, while $\ell=4$ and every $\ell\ge6$ can.
\end{proposition}
\begin{proof}
Suppose a component of length $5$ occurred; then only positions $p$ and
$p+1$ fired in it, both at least once, and no other position within
distance $4$ fired (any such firing would lengthen the component). Note
that window $p+1$ spans sites $[p{+}1,p{+}4]$, all inside the component,
and the only positions touching those sites are $p{-}2,p{-}1,p,p{+}1$; the
first two are excluded by the same argument, so only $p$ and $p+1$ can
alter the activity of $p+1$. Consider the first firing among $\{p,p{+}1\}$,
say $p$ (the other case is symmetric using (iii)). At that instant $p$ is
active, so by (i) $p+1$ is not. After the firing, by (ii), $p+1$ is still
inactive. Further firings of $p$ only toggle its window between $A$ and $B$,
and (ii) applies to each; $p+1$ therefore never becomes active and never
fires, contradicting the assumption. Lengths $4$ and $\ge6$ are realisable
because a single position, or positions at offset $2,3,4,\dots$, can fire.
\end{proof}

This reproduces the observed support exactly: $\{4\}\cup\{6,7,8,\dots\}$. It
is consistent with, and sharper than, the thesis result
$\ell_{\min}=\ell_{\max}=4$ for the minimal two-particle bubble, which
explains $\ell\le3$ and the prominence of $\ell=4$ but does not address
$\ell=5$.

\medskip
\noindent\fbox{\parbox{0.97\textwidth}{\small
\textbf{Status of Proposition~\ref{prop:five}.} The exhaustive verification
of Lemma~\ref{lem:offset1} is machine-checked and reproducible. The
combinatorial argument built on it was written while preparing this
document and \emph{has not been reviewed by the researcher}. Per the
project's own rule that theoretical arguments are confirmed before being
relied on, it should be checked before being cited or built upon.}}

\section{Claims, graded by evidence}

\subsection*{Established --- directly supported by data in this repository}
\begin{itemize}
\item \Est\ \textbf{Logarithmic decay (H1c) is excluded.}
$\chi^2_\nu=2776$ (definition A), $420$ and $543$ (definition B).
\emph{Evidence:} Table~\ref{tab:h1fits}, Fig.~\ref{fig:h1models}.
\item \Est\ \textbf{Any power-law decay of $f_N$ is bounded to
$\alpha<1.5\times10^{-4}$ at 95\% confidence over $N=10^2$--$3\times10^5$.}
\emph{Evidence:} H1b fit, definition A.
\item \Est\ \textbf{The frozen fraction is $0.359$--$0.363$ depending on
definition, with a definitional spread of ${\sim}0.004$.}
\emph{Evidence:} three independent weighted means; the ordering
$f_A<f_B(10L)<f_B(5L)$ obeys the required monotonicity.
\item \Est\ \textbf{$f_B(W)$ saturates for $W\gtrsim20L$;} the production
$W=5L$ overshoots by $0.003$--$0.004$.
\emph{Evidence:} Table~\ref{tab:wdep}, Fig.~\ref{fig:wdep}.
\item \Est\ \textbf{The dynamical estimator reproduces exact static
reachability} on the sub-ensemble it can sample, $\chi^2_\nu=0.83$ over six
$L$. \emph{Evidence:} Table~\ref{tab:validation}.
\item \Est\ \textbf{The model is strongly fragmented:} sector/symmetry-sector
$\sim(0.576)^L$. \emph{Evidence:} Fig.~\ref{fig:bfs}(c).
\item \Est\ \textbf{Active-region lengths $1,2,3,5$ cannot occur.}
\emph{Evidence:} zero counts in ${\sim}57{,}000$ pooled gaps at eight $L$,
plus the exhaustive Lemma~\ref{lem:offset1} (the argument assembling them,
Proposition~\ref{prop:five}, is pending review).
\item \Est\ \textbf{An exponential describes $P_L(\ell)$ better than a power
law over the full range, at every $L$.} \emph{Evidence:}
Table~\ref{tab:h2}, Fig.~\ref{fig:h2pmf}.
\end{itemize}

\subsection*{Strong observations --- robust but needing more verification}
\begin{itemize}
\item \Obs\ \textbf{$f_N$ saturates to a nonzero constant (H1a).} Flat over
more than three decades, under two independent definitions, with an
exclusion of the alternatives. What keeps this from ``established'' is that
no measurement of a decay slower than $\alpha\approx10^{-4}$ is possible
with this range, and the underlying quantity is only \emph{known} to be
definition-independent up to the $0.004$ spread.
\item \Obs\ \textbf{The active-region length distribution is $L$-independent
in shape (H2a over H2b).} Collapse in Fig.~\ref{fig:h2collapse}; $\lambda$,
$\langle\ell\rangle$ and $P(\ell{=}0)$ flat over two decades. Weakened by
the fact that the far tail alone does not discriminate, and the largest $L$
rests on only 20 realisations.
\item \Obs\ \textbf{The settings shift is entirely a $W$ effect.} The
$W$-scan accounts for $0.003$--$0.004$ of the observed $0.00298$; nothing is
left for a burn-in explanation.
\end{itemize}

\subsection*{Conjectures / interpretations not established}
\begin{itemize}
\item \Conj\ \textbf{$f_\infty\simeq0.36$ is the thermodynamic-limit static
frozen fraction.} The exact calculation only reaches $L=60$, where
$f_{\text{static}}=0.337\pm0.010$ and still rising; the identification of
the measured plateau with the $L\to\infty$ static value is an extrapolation.
\item \Conj\ \textbf{Because active regions are exponentially distributed
and $L$-independent, global transport is arrested with no anomalous
hydrodynamic tail.} This is the README's motivating inference, and the data
support the premise --- but the conclusion concerns density--density
correlation functions, which have never been computed in this project.
\item \Conj\ \textbf{The measured exponential is the thesis's
$\nu^x(1-\nu)^{\ell-x}$ law.} Plausible and the forms match, but the
thesis's prediction has never been evaluated numerically at $\nu=1/2$ and
compared point-by-point against our $P_L(\ell)$.
\item \Conj\ \textbf{$\ell=5$ is forbidden at every particle number $x$.}
Proposition~\ref{prop:five} proves it for the dynamical active runs we
measure; the literature note records a separate hand-argument for $x=3$
only, which was never verified against our data.
\end{itemize}

\subsection*{Claims with currently insufficient evidence}
\begin{itemize}
\item Any statement about $N>10^4$ under the current definition. The largest
definition-B run is $N=10^4$; everything beyond that is definition A.
\item Any statement about densities other than $1/2$. Nothing else has been
run, yet the literature's whole framework is organised around filling.
\item The precise value ``$0.36\pm0.009$'' recorded in
\texttt{HYPOTHESES.md}. The $\pm0.009$ has no stated derivation; the
inverse-variance weighted errors are $\pm0.0007$, and the definitional
spread is $\pm0.004$. The quoted uncertainty appears to be a conservative
eyeball rather than a computed quantity.
\end{itemize}

\section{The current scientific picture}

\paragraph{Question.} Does a dipole-conserving hardcore chain at half
filling freeze a finite fraction of its sites in the thermodynamic limit,
and are the surviving active regions microscopic or scale-free?

\paragraph{Approach.} Direct stochastic simulation of the constrained
dynamics, with the frozen set identified by a trailing-window criterion
after a self-scaling burn-in; ensembles sized adaptively to a 1\% relative
error; and, at small $L$ where it is possible, exact enumeration of the
reachable configuration space as ground truth.

\paragraph{Results.} $f_N$ is flat at $0.359$--$0.363$ from $N=10^2$ to
$3\times10^5$. Logarithmic decay is excluded; power-law decay is bounded
below $10^{-4}$ in exponent. The estimator is validated against exact
enumeration. Active-region lengths are exponentially distributed with
$\lambda\approx13$ sites, $L$-independent, supported on
$\{4\}\cup\{6,7,\dots\}$.

\paragraph{Observation.} Two independent structural facts coexist: a finite
frozen backbone occupying ${\sim}36\%$ of the chain, and active clusters
whose size distribution has a fixed microscopic scale that does not grow
with the system. Exact enumeration shows the accessible configuration space
is an exponentially small fraction of the symmetry sector.

\paragraph{Interpretation (labelled as such).} The chain is strongly
Hilbert-space fragmented at half filling. Because the active regions have a
finite correlation length, the system does not possess the rare large
clusters that would be needed to carry long-range relaxation through a
frozen background. In the README's framing this is the ``exponential''
branch: transport is arrested, and no anomalous tail in the global
density--density correlator is expected \emph{from this mechanism}.

\paragraph{Claim.} H1a (saturation) is supported and its competitors are
excluded within the measured range; H2a (an $L$-independent, memoryless
active-region distribution) is supported over H2b. The forbidden-length
structure is a consequence of the 4-site move geometry rather than a
statistical accident.

\section{Contradictions, weaknesses and problems}
\label{sec:problems}

This section is deliberately unflattering.

\begin{enumerate}
\item \textbf{\texttt{HYPOTHESES.md} H1 still defines the wrong
quantity.} Its order-parameter paragraph defines $f_N$ statically
(``no sequence of allowed moves can ever change its occupation \dots a
property of the initial string alone, not of a stochastic trajectory or run
time''), but every production number in the project is a dynamical proxy
with an explicit run-time parameter. Section~\ref{sec:bfs} now shows the two
agree, but the document has never been corrected, and a reader taking H1 at
face value would misunderstand what was measured. The \Unc\ block in H1 is
also now out of date.

\item \textbf{The ``$0.36\pm0.009$'' in \texttt{HYPOTHESES.md} is not a
computed error bar.} Nothing in the repository derives $\pm0.009$. The
statistical error is $\pm0.0007$ and the definitional spread $\pm0.004$.
Carrying an unexplained uncertainty into a recorded claim is a hazard if it
is later quoted.

\item \textbf{The settings comparison was a confounded experiment.} Four
parameters were changed simultaneously; two of them independently predict
the observed sign. It is resolved only retrospectively, by the $W$-scan.
A one-at-a-time design was proposed at the time and not run.

\item \textbf{Heterogeneous precision is silently pooled.} The 24-point
``definition B'' dataset mixes \texttt{test.yaml} points (10\% target
relative error) with production points (1\%). The two visible outliers at
$N=200,400$ in Fig.~\ref{fig:h1overview} are exactly those. Inverse-variance
weighting limits the damage, but the exploratory notebook's headline
``weighted mean over 24 points'' is not a homogeneous dataset.

\item \textbf{Definition A's constant fit is formally a poor fit.}
$\chi^2_\nu=1.57$ on 30 dof is $p=0.024$. Either the error bars are
underestimated --- plausible, since each point's error is a
realisation-to-realisation standard error that ignores any systematic
residual burn-in bias --- or there is real residual structure. The project
investigated the latter (commensuration, second correction term,
log-periodicity) and found nothing statistically supported, which leaves
mildly underestimated uncertainties as the most likely explanation. That
possibility has never been tested directly.

\item \textbf{Sampling at the largest sizes is thin.} At $L=12800$, H2 rests
on 20 realisations; the H1 $W=5L$ run has $n=20$ at eleven of twenty
$N$-values. Adaptive sampling stops at a target \emph{relative} error of the
mean, which at large $L$ is met almost immediately because
realisation-to-realisation variance shrinks --- so the largest, most
physically interesting sizes get the \emph{fewest} samples. This is a
structural feature of the stopping rule worth reconsidering.

\item \textbf{A real discard bias exists and is only extrapolated away at
production sizes.} Jammed configurations are discarded. This is a
$4.7\sigma$ effect at $L=16$. It is believed negligible at $L\ge100$
because \texttt{n\_discarded}$=0$ in every production run and the jamming
probability decays exponentially --- but ``zero discards observed'' with
$n\le1400$ only bounds the jammed fraction at ${\lesssim}10^{-3}$, it does
not measure it.

\item \textbf{The $W\to\infty$ plateau is a plateau over one decade only.}
Flatness of $f_B(W)$ for $W\in[20L,200L]$ does not prove convergence to the
static value; a site with a flip timescale $\gg200L$ would be invisible. The
small-$L$ agreement with exact enumeration is reassuring, but at $L\le56$,
where absolute timescales are tiny.

\item \textbf{The tail-only fit does not support the headline H2 claim, and
this asymmetry is under-emphasised.} Restricted to $\ell\ge15$, the power
law is marginally \emph{better} at five of eight $L$. The full-range
comparison is more discriminating for good reasons, but the honest
statement is ``exponential over the full range; the far tail alone is
inconclusive'', not ``exponential, confirmed''. No AIC or likelihood-ratio
test has been performed on the raw samples.

\item \textbf{The archived analysis notebooks are dead code.} The three
\texttt{h1\_fit\_frozen\_fraction*.ipynb} notebooks are pinned to paths that
have moved to \texttt{data/archive/}; they cannot be re-executed. The
definition-A fits quoted throughout this document were therefore recomputed
from the raw \texttt{.npy} files rather than trusted from notebook output.
(They reproduce the logged values.)

\item \textbf{Stale statements in the research log.} The entry of 13 Sep
states that \texttt{production\_large\_scale.yaml} is ``scheduled to launch
automatically at 2026-09-14 00:00''. It never ran, and the config has since
been retargeted twice. Log entries are a historical record, so this is not
wrong in itself, but nothing in the log records the interruption.

\item \textbf{The core dynamics had never been tested, and there is still
no regression test.} Every number in the project flows through one
\texttt{DipoleChain} class, whose conservation laws were asserted only in
its docstring. I ran that check ad hoc while preparing this document and it
passes cleanly (Section~\ref{sec:results}, check~1), and the independently
written BFS code agrees on the move rule --- so this is no longer an open
risk. But the check exists only in this document's history: there is no test
file in the repository, so nothing will catch a future regression.

\item \textbf{The project's stated physical target has not been approached.}
The README's actual goal is the density--density correlation function and
its possible anomalous tail. Two structural inputs have now been measured,
but no correlation function, no transport coefficient and no dynamical
exponent has been computed.
\end{enumerate}

\section{Highest-value next steps, in priority order}

\begin{enumerate}
\item \textbf{Commit the conservation-law check as a regression test.}
(Minutes.) The check described in Section~\ref{sec:results} passes, but it
lives nowhere in the repository. Turning it into a test file --- assert
$\sum_i s_i$ and $\sum_i i\,n_i \bmod L$ invariant under \texttt{step()},
and the incremental active-list update equal to a full rescan --- costs
almost nothing and protects every number the project produces. Highest
value-per-cost in the list.

\item \textbf{Run \texttt{production\_large\_scale.yaml}.} (4--5\,h,
already costed and approved.) It is the only planned measurement that
extends definition B beyond $N=10^4$, and the H1 claim currently has a
one-decade gap between where the current definition stops ($10^4$) and where
the superseded one reached ($3\times10^5$).

\item \textbf{Re-weight the sampling toward large $L$.} Replace, or
supplement, the relative-SE stopping rule with a minimum sample count that
\emph{grows} with $L$. At present the largest systems are the least
sampled, which is backwards for a finite-size study.

\item \textbf{Do proper model selection on H2 with likelihoods.} Fit the
raw $\ell$ samples by maximum likelihood to a geometric/exponential and to
a discrete power law (with the $\ell\in\{4\}\cup\{6,\dots\}$ support handled
explicitly), and compare with AIC and a likelihood-ratio test. The current
$R^2$-on-binned-counts comparison is weaker than the data deserve and is
sensitive to the binning and to the single-count floor.

\item \textbf{Evaluate the thesis's exact prediction at $\nu=1/2$ and
overlay it on Fig.~\ref{fig:h2pmf}.} The formula
$\propto\nu^x(1-\nu)^{\ell-x}$ plus $P_b(\nu)$ is closed-form; computing it
and plotting it against our measured $P_L(\ell)$ converts a qualitative
``consistent with the literature'' into a quantitative test with no new
simulation.

\item \textbf{Push the exact enumeration further with symmetry reduction.}
$L=60$ took 6.7 minutes; using translation symmetry and canonical forms
should reach $L\approx80$--$100$ and would let $f_{\text{static}}(L)$ be
extrapolated against the dynamical plateau directly rather than
qualitatively.

\item \textbf{Measure the density--density correlation function.} This is
the README's actual target and the only way to convert the structural
results into the physical claim the project wants to make.

\item \textbf{Vary the density.} Everything is at $\nu=1/2$. The literature
frames this model entirely in terms of filling, and the thesis's claim that
$(d,k)=(2,4)$ is strongly fragmented \emph{at every} filling is directly
testable with the existing pipeline by changing one config value.
\end{enumerate}

\section{Questions I could not resolve from the project}

\begin{enumerate}
\item \textbf{Which frozen-fraction number is the one you want to quote?}
There are now three defensible candidates: $0.35988\pm0.00015$
(definition A, widest $N$ range), $0.36303\pm0.00065$ (definition B at the
production $W=5L$), and the $W$-saturated plateau $\approx0.3647$ at
$L=4000$. They differ by more than their statistical errors. I have not
picked one as authoritative anywhere in this document.

\item \textbf{Where does the $\pm0.009$ in \texttt{HYPOTHESES.md} come
from?} I could not reconstruct it from any file. If it was meant as a
conservative band covering the definitional spread, it should say so; if it
was meant as an error bar, it disagrees with the computed one by an order of
magnitude.

\item \textbf{Should ``frozen'' ultimately mean static or rolling?} They now
agree numerically, so this is a presentational choice --- but it changes what
H1 claims. The static definition connects directly to the
Hilbert-space-fragmentation literature; the rolling one is what is actually
measurable at large $L$.

\item \textbf{Is discarding zero-length gaps the right convention for the
physics you care about?} It removes ${\sim}89\%$ of all gaps and breaks the
$\langle\ell\rangle=(1-f)/f$ identity. For a correlation-function
calculation the full distribution including zeros is what enters. Was the
intent to characterise the shape of nonzero clusters only?

\item \textbf{Is the README's ``aperiodic'' a typo for ``a periodic''?} The
code is unambiguously periodic, and everything in the project assumes PBC,
so I have read it as a typo --- but if an aperiodic or quasiperiodic variant
was ever intended, none of the current work addresses it.

\item \textbf{Do you want Proposition~\ref{prop:five} checked
independently before it is used?} Per your own rule I have flagged it rather
than treating it as established.
\end{enumerate}

\section*{Your suggestions / things you should consider}
\addcontentsline{toc}{section}{Your suggestions / things I should consider}

\emph{These are suggestions, not results.}

\begin{itemize}
\item \textbf{Measure the distribution of frozen-run lengths too, not just
active-run lengths.} You currently characterise the gaps but not the
backbone. The frozen runs are the complementary structure, they are trivially
available from the same snapshot (the code already builds the mask), and
their distribution is what determines whether the frozen backbone is
percolating-like or fragmented. Given $P(\ell=0)\approx0.888$ the frozen
runs are clearly long, and nobody has looked at them.

\item \textbf{The zero-gap rate is suspiciously constant.} $0.8850$ to
$0.8888$ across two decades of $L$, i.e.\ variation in the fourth decimal.
Combined with $\langle\ell\rangle$ and $\lambda$ also being flat, this
smells like an exactly computable constant of the model rather than an
empirical number. It may be derivable from the same window-geometry
argument as Proposition~\ref{prop:five}.

\item \textbf{Exploit the ``all $W$ from one trajectory'' trick more
widely.} It converted a noisy 4-parameter comparison into a clean paired
measurement at a fraction of the cost. The same pattern --- record per-site
last-flip times once, derive many observables --- would give you the
autocorrelation function $C(t)$ and the distribution of single-site flip
waiting times essentially for free from runs you are already doing. The
waiting-time distribution in particular would tell you whether the
$W$-plateau is genuine or hiding a slow tail.

\item \textbf{Consider reporting $1-f$ rather than $f$.} The active
fraction $0.64$ is the quantity that enters transport arguments, and error
propagation is more transparent there.

\item \textbf{A theoretical connection worth checking.} Your active regions
are bounded by frozen sites and have a minimum size of 4 with $5$ forbidden.
That is the signature of a constrained-tiling problem. The Moudgalya--Motrunich
technique in your notes (canonical form $\to$ tiling $\to$ transfer matrix)
would plausibly yield $P(\ell)$ in closed form for this model, including the
exact $\lambda$, and you already have the numerical answer to check it
against. This looks like the single most tractable piece of analytics
available to you right now.

\item \textbf{Be careful before abandoning the model.} The README says you
would move on if the distribution is exponential. It is --- but the
exponential decay length $\lambda\approx13$ sites is large compared to the
4-site move, and $89\%$ of frozen sites are adjacent to another. Before
concluding the model is uninteresting for hydrodynamics, it is worth
checking whether the \emph{frozen} structure has long-range correlations
even though the active clusters do not.

\item \textbf{Housekeeping.} The archived notebooks cannot be re-executed;
either repoint them at \texttt{data/archive/} or add a line to their README
stating they are frozen artifacts. And the two configured-but-never-run
jobs should either be run or removed from the configs directory, since a
config that has never executed is a trap for a future reader.
\end{itemize}

\vfill
\noindent\rule{\textwidth}{0.4pt}\\
{\footnotesize
Figures regenerable via
\texttt{python numerics/scripts/make\_progress\_figures.py}.
Numerical values dumped to
\texttt{numerics/data/processed/progress\_numbers.json}.
New measurements made for this document live under
\texttt{numerics/data/raw/} in the
\texttt{static\_frozen\_*} and \texttt{w\_dependence\_*} run directories,
and are logged in \texttt{RESEARCH\_LOG.md} (14 September entries).
Document source: \texttt{progress.tex}.}

\end{document}
